\subsection{应用于量子化}\label{qtz}

接下来我们想得到与第 \ref{symmeq} 节形变量子化结果类似的量子化结果。为此，我们特殊化到这样的情形：$Y$ 的复辛同胚 $\varphi$ 是 $M$ 的实辛同胚 $\varphi_r$ 的解析延拓。我们假设 $\varphi_r$ 提升为预量子线丛 $\frL\to M$ 的对称；换言之，假设 $\varphi_r^*(\frL)$ 与 $\frL$ 作为带联络的酉线丛同构。记 $L_r$ 为 $M\times M'$ 中与 $\varphi_r$ 相伴的拉格朗日对应，$L$ 为其在 $Y\times Y'$ 中的复化。设 $\I_\varphi$ 是 $Y$ 的 $A$-模型中相应的拓扑界面。上一节讨论了 $\I_\varphi$ 在 $\B_\cc$ 上的端点 $\alpha$ 的性质，使 $\I_\varphi$ 实现代数 $\A$ 的自同构。为了使 $\I_\varphi$ 不仅作用于 $\A$，而且作用于量子 Hilbert 空间，即作用于 $A$-模 $\H=\Hom(\B,\B_\cc)$，我们还需要一个端点 $\beta$，使 $\I_\varphi$ 在支撑于 $M$ 的拉格朗日膜 $\B$ 上结束。

在折叠图像中，$\B$ 来自支撑在 $M\times M'\subset Y\times Y'$ 上的拉格朗日膜 $\h\B$，其 Chan--Paton 丛为 $\frL^{-1}\boxtimes\frL$（第二个因子是第一个的逆，因为折叠把 $\CP$ 丛换成其对偶）。界面 $\I_\varphi$ 在折叠图像中来自支撑在 $L$ 上的拉格朗日膜 $\B_L$，其 $\CP$ 丛平凡。角空间 $\Hom(\B_L,\h\B)$ 的首阶近似，就是交 $L\cap(M\times M')$ 上取值于 $\CP$ 丛张量积的 de Rham 上同调。由于 $\B_L$ 的 $\CP$ 丛平凡，相关的 $\CP$ 丛之积就是 $\frL^{-1}\boxtimes\frL$。

交 $L\cap(M\times M')$ 就是实拉格朗日对应 $L_r$。由于 $L_r$ 参数化 $M\times M'$ 中形如 $(m,\varphi_r(m))$ 的对，它在拓扑上是 $M$ 的一个拷贝。为简单起见，假设 $b_1(M)=0$（我们继续假设 $b_1(Y)=0$）。这意味着 $b_1(L_r)=0$，且 $A$-模型空间 $\Hom(\B_L,\h\B)$ 就是 $L_r$ 上取值于 $\frL^{-1}\boxtimes\frL$ 的 de Rham 上同调。作为 $M\times M'$ 上的线丛，$\frL^{-1}\boxtimes\frL$ 当然不平凡；但 $\varphi_r^*(\frL)$ 与 $\frL$ 同构的假设意味着，限制到 $L_r$ 上时 $\frL^{-1}\boxtimes\frL$ 平凡。$\varphi_r^*(\frL)$ 与 $\frL$ 之间的同构在乘以模 1 复数的意义下唯一。每个同构 $\beta:\frL\to\varphi^*(\frL)$ 的选择，都确定 $H^0_{\mathrm{dR}}(L_r,\frL^{-1}\boxtimes\frL)$ 的一个元素 $\beta$。这些就是 $\I_\varphi$ 与 $\B$ 之间的典则端点。

这类角在 $\sigma$-模型映射 $\Phi:\Sigma\to Y$ 的语言中有简单含义。$\Phi$ 把二维流形 $\Sigma$ 映到 $Y$，同时把由 $\B$ 标记的 $\partial\Sigma$ 的一个分支映到 $M$。在界面 $\I_\varphi$ 与 $\partial\Sigma$ 相遇的端点 $\sigma$ 处，映射 $\Phi:\partial\Sigma\to M$ 不连续，从 $m$ 跳到 $\varphi(m)$。重要的是，若用类 $\beta\in H^0_{\mathrm{dR}}(L_r,\frL^{-1}\boxtimes\frL)$ 定义端点，则对 $\Phi(\sigma)$ 的行为没有进一步限制。反之，若用例如 $H^k_{\mathrm{dR}}(L_r,\frL^{-1}\boxtimes\frL)$ 中 Poincaré 对偶于余维 $k$ 子流形 $V\subset L_r$ 的元素，则我们必须施加约束 $\Phi(\sigma)\in V$。

“无约束”条件在合成界面时能很好地自我复制。设 $\varphi_r,\varphi_r'$ 是 $M$ 的两个实辛同胚，都提升为 $\frL$ 的对称；定义 $\varphi_r''=\varphi_r'\circ\varphi_r$。在体内，与 $B$-模型的关系表明相应界面按自然方式合成：$\I_{\varphi'}\cdot\I_\varphi=\I_{\varphi''}$。我们希望给它们配上与 $\B$ 的端点，使这一简单合成律在界面于 $\B_\cc$ 上结束时仍成立。为此，我们使用与同构 $\beta:\frL\to\varphi_r^*(\frL)$、$\beta':\frL\to\varphi_r'{}^*(\frL)$ 相伴的简单端点。如上段所述，当界面 $\I_\varphi$ 或 $\I_{\varphi'}$ 在 $\partial\Sigma$ 上结束时，这样的端点不施加约束。“无约束”与“无约束”的合成仍是“无约束”。若用 $\beta''=\beta'\circ\beta$ 定义 $\beta'':\frL\to\varphi_r''{}^*(\frL)$，则包括其在 $\partial\Sigma$ 上的端点在内，界面的合成就简单地是 $\I_{\varphi'}(\beta')\cdot\I_\varphi(\beta)=\I_{\varphi''}(\beta'')$。

关于这个答案，有几点值得强调。第一，对这些界面在拉格朗日膜 $\B$ 上的端点，这个答案比我们在其于 $\B_\cc$ 上的端点所得到的答案直接得多。第二，为得到这个答案，我们必须考虑由辛同胚 $\varphi_r:M\to M$ 连同同构 $\beta:\frL\to\varphi^*(\frL)$ 组成的对。这样的对正是经典力学对称群 $\G$ 的元素，如第 \ref{problem} 节所定义。

现在设 $\T$ 是 $Y$ 的复辛同胚的集合，这些辛同胚在配上与 $\B_\cc$ 的合适角后生成 $\A$ 的自同构。于是对任何 $\varphi\in\T$，界面 $\I_\varphi$ 有端点 $\alpha$，使 $\I_\varphi(\alpha)$ 生成自同构。我们在第 \ref{symmeq} 节得知 $\T$ 是群，而且作为 $\A$ 的自同构群作用。设 $\T_r$ 是 $\T$ 中这样的辛同胚构成的子群：它们把 $M$ 映到自身，且限制到 $M$ 上时提升为 $\frL$ 的对称。由于与辛同胚相伴的 $\frL$ 对称只在整体相位意义下唯一，$\T_r$ 出现在中心扩张
\be\label{zirko}1\to\UU(1)\to\T_r^*\to\T_r\to1\ee
中。每个 $\varphi\in\T_r$ 都是某个 $\varphi_r:M\to M$ 的解析延拓，使得存在同构 $\beta:\frL\to\varphi_r^*(\frL)$。这个同构可用于定义 $\I_\varphi$ 在 $\B$ 上的端点，我们也称之为 $\beta$。给每个 $\varphi\in\T_r$ 的 $\I_\varphi$ 配上这样的端点 $\alpha,\beta$，我们定义作用在 $\H=\Hom(\B,\B_\cc)$ 上的量子算符 $\O_{\varphi,\alpha,\beta}$。鉴于关于体算符与端点的讨论，这些算符 $\O_{\varphi,\alpha,\beta}$ 能正确合成，给出群 $\h\T_r$ 在 $\H$ 上的作用。在第 \ref{symmeq} 节末尾，我们对 $M=\R^{2n}$ 演示了这一点，其方式依赖于该情形中代数 $\A$ 相当大这一事实。然而，即使 $\A$ 不那么大，结果仍成立。
